\begin{figure}[H]
        \centering

        \begin{subfigure}[t]{0.3\textwidth}
            \centering
            \includesvg[width=\textwidth]{img/graphs/6/F2/1}
            \caption{
                Plocha S = (140,7 ± 145,8) $\mathrm{mN \cdot s}$
            }
            \label{fig:integrales-2a}
        \end{subfigure}%
    \hfill
\begin{subfigure}[t]{0.3\textwidth}
            \centering
            \includesvg[width=\textwidth]{img/graphs/6/F2/2}
            \caption{
                Plocha S = (138,1 ± 145,9) $\mathrm{mN \cdot s}$
            }
            \label{fig:integrales-2b}
        \end{subfigure}%
    \hfill
\begin{subfigure}[t]{0.3\textwidth}
            \centering
            \includesvg[width=\textwidth]{img/graphs/6/F2/3}
            \caption{
                Plocha S = (137,8 ± 142,6) $\mathrm{mN \cdot s}$
            }
            \label{fig:integrales-2c}
        \end{subfigure}

    \vspace{0.5em}
\begin{subfigure}[t]{0.3\textwidth}
            \centering
            \includesvg[width=\textwidth]{img/graphs/6/F2/4}
            \caption{
                Plocha S = (135,7 ± 136,7) $\mathrm{mN \cdot s}$
            }
            \label{fig:integrales-2d}
        \end{subfigure}%
    \hfill
\begin{subfigure}[t]{0.3\textwidth}
            \centering
            \includesvg[width=\textwidth]{img/graphs/6/F2/5}
            \caption{
                Plocha S = (134,5 ± 129,9) $\mathrm{mN \cdot s}$
            }
            \label{fig:integrales-2e}
        \end{subfigure}%
    \hfill
\begin{subfigure}[t]{0.3\textwidth}
            \centering
            \includesvg[width=\textwidth]{img/graphs/6/F2/6}
            \caption{
                Plocha S = (137,3 ± 142,5) $\mathrm{mN \cdot s}$
            }
            \label{fig:integrales-2f}
        \end{subfigure}

    \vspace{0.5em}
\begin{subfigure}[t]{0.3\textwidth}
            \centering
            \includesvg[width=\textwidth]{img/graphs/6/F2/7}
            \caption{
                Plocha S = (163,8 ± 155,8) $\mathrm{mN \cdot s}$
            }
            \label{fig:integrales-2g}
        \end{subfigure}%
    \hfill
\begin{subfigure}[t]{0.3\textwidth}
            \centering
            \includesvg[width=\textwidth]{img/graphs/6/F2/8}
            \caption{
                Plocha S = (138,2 ± 137,3) $\mathrm{mN \cdot s}$
            }
            \label{fig:integrales-2h}
        \end{subfigure}%
    \hfill
\begin{subfigure}[t]{0.3\textwidth}
            \centering
            \includesvg[width=\textwidth]{img/graphs/6/F2/9}
            \caption{
                Plocha S = (130,9 ± 135,0) $\mathrm{mN \cdot s}$
            }
            \label{fig:integrales-2i}
        \end{subfigure}

    \vspace{0.5em}
\begin{subfigure}[t]{0.3\textwidth}
            \centering
            \includesvg[width=\textwidth]{img/graphs/6/F2/10}
            \caption{
                Plocha S = (137,3 ± 135,1) $\mathrm{mN \cdot s}$
            }
            \label{fig:integrales-2j}
        \end{subfigure}

        \caption{
            Naměřené hodnoty síly $F$ v závislosti na čase $t$. Výpočet plochy pod integrálem. Vypočteno skrze program Python \cite{GitHubMAT}.
        }
        \label{fig:integrales-2}
    \end{figure}
    